\section{The two ambient lattices and the construction of a good core}

% BEGIN P78C1
\paragraph{P78C1. Recovering a full lattice from the given complex reading.}
Let $L$ be the coefficient number field, $\sigma:L\hookrightarrow
\mathbb C$ the specified embedding and $j:L^n\to W$ a
$\sigma$-semilinear map. Suppose the $j(e_i)$ are independent over
$\mathbb C$. Let $S$ be the finite generating set of the given
ambient lattice, spanning $W$ over $\mathbb C$. The source
lattice comparison says that every $s\in S$ belongs to the range
of $j$ and, for the specified submodule $N\subset L^n$,
\[
 x\in N \quad\Longleftrightarrow\quad
 j(x)=\sum_{s\in S}\sigma(a_s)s
 \quad\hbox{for some }a_s\in\mathcal O_L .
\]
Then $N$ is finitely generated over $\mathcal O_L$ and spans $L^n$.
Indeed $j(x)=\sum_i\sigma(x_i)j(e_i)$ is injective. The unique
preimages $u_s=j^{-1}(s)$ give
$N=\sum_s\mathcal O_L u_s$. If their $L$-span were proper,
choose a nonzero $L$-linear functional $d$ killing it. On the
complex span of the $j(e_i)$ define $d_{\mathbb C}$ by the values
$\sigma(d(e_i))$ on these independent vectors, then extend it
linearly to $W$. Thus $d_{\mathbb C}(j(x))=\sigma(d(x))$.
It vanishes on $S$, hence on $W$, forcing $d=0$. The same
argument covers a zero-dimensional ambient space.
% END P78C1

% BEGIN P78C2
\paragraph{P78C2. The metric at every positive place and the norm at the reading place.}
Let $L$ be totally complex. For the chosen matrix $H\in M_n(L)$
and an infinite place $v$ let $\tau$ be its representative embedding
and put $h_v(x,y)=\overline{\tau(x)}^{\,t}\tau(H)\tau(y)$.
The source tests $\mathrm{PosAt}(H,v)$ on $x,y\in L^n$.
If it holds, testing the coordinate vectors shows that $\tau(H)$
is Hermitian. Since $\tau(L)$ contains $\mathbb Q$ and a nonreal
element, it is dense in $\mathbb C$. Continuity thus gives
positive semidefiniteness on $\mathbb C^n$. If $\tau(H)$ were
singular, $\det H=0$ and there would be a nonzero rational vector
$x\in L^n$ with $Hx=0$, contradicting strict positivity on
that vector. Thus $\tau(H)$ is positive definite.

Use $X_v=\tau(H)^{-1/2}$ at places with positive definite
matrix and $X_v=1$ elsewhere. These frames are invertible.
The majorant has Gram matrix $(X_v^{-1})^* X_v^{-1}$, hence
equals $h_v$ at every place satisfying the source positivity test.
Suppose $(j,H)$ has the prescribed inner-product comparison at
$\sigma$, which sends the field involution to complex conjugation.
The representative of $v_\sigma$ is either $\sigma$ or
$\bar\sigma$. Thus $h_{v_\sigma}(x,y)$ is the inner product
of $j(x),j(y)$ or its conjugate, with the same choice for all
$x,y$. Its diagonal is $\|j(x)\|^2$, strictly positive for
$x\ne0$ by injectivity in P78C1. Hermitian symmetry also follows.
Taking nonnegative square roots in the resulting majorant identity
gives $\|x\|_{v_\sigma}=\|j(x)\|$. No metric comparison at
a place lacking $\mathrm{PosAt}$ is needed.
% END P78C2

% BEGIN P78C3
\paragraph{P78C3. Both ambient forms for the actual row condition.}
Fix a fourth-stage record $l$ on a pinned column, with row slot
$\mathcal S.R_x=\mathrm{RX}$. For each full tuple $\mu$ and each
of the two indicated levels, apply P78C1 to
$N=\mathrm{rNab}$ and to the all-tuple reading stored in $l.rx$.
Its form tie supplies the independence and complex spanning of the
coordinate vectors. Its lattice tie supplies the finite generator
comparison. Equip $L^{\mathrm{rDim}}$ with that lattice, its
coordinate basis and the frames of P78C2 for $H=\mathrm{rH}$.
The embedded map is the very reading $j$ from the row. Finite
generation, full span and invertibility give a good Hermitian
lattice. The remaining ambient-form fields are the basis, generator,
membership and norm identities proved above or supplied by the tie.

Restrict this construction to the tuples used in the row packet.
For any such tuple $s$, the inclusion $s\mapsto s$ into all tuples
agrees with the full tuple of $Y.\mathrm{tup}(s)$: the index
identity of the reading picks a representative whose tuple is $s$.
Thus the coordinate dimensions agree, as do $\mathrm{rNab}$ and
$\mathrm{rH}$ after reindexing. Take the coordinate equivalence
which sends entry $i$ to entry $i$ under this dimension equality.
The restricted all-tuple map agrees with the installed map by
$\mathrm{SameOn}$, precisely agreement of coordinates with the
same integer index. This proves the equality of readings. The
submodule transports by this equivalence and the majorants agree
on both arguments at every positive place by P78C2. These are
the three fields of $\mathrm{AmbTie}$.

At level one take the sharp row side, its subgroup and character,
and $l.rx.\mathrm{all}.j_A$. At level zero take $l.p_0.s$,
its subgroup and character and $l.rx.\mathrm{all}.j_{A,0}$.
The two restrictions use the two separate map-agreement fields
of $l.rx$. This constructs $a,a_0$ with $\mathrm{AmbAre}$ for the
original lift reading. It proves the complete $\mathrm{AmbEx}$
clause for any such slot system, in particular for both passes
of the $\mathrm{prmR}$ construction. The construction works on
every admitted fourth-stage record and makes no field enlargement
or further restriction on the size parameter.
% END P78C3

% BEGIN P78C4
\paragraph{P78C4. Vanishing off the selected places in the fourth-stage core.}
Write $D=w.D(Z)$ and $t=l.\mathrm{what}_2$. For every size record
$L_0$ of scale $s=K.s$, every $v\notin l.W_t$ and
$1\le j\le D.R.m$ one has $\mathrm{Fval}(L_0,j,v)=0$.
Indeed let $p$ be the residue characteristic of $v$. The formula
for forced content first tests whether $p$ is one of the two
distinguished type primes. This is exactly $p=g(i)$ for a binary
index $i$. The remaining disjoint branches in the ordered formula
give $F_{\rm RamD}$, the bounded inert term $F_{\rm Inert}$, or
zero. Hence $\mathrm{Fval}(L_0,j,v)=\mathrm{FK}(s,D,j,p)$.
The row gives the equality $l.W_t=z.W_1(F_{0,K})$. In this set's
definition take the prime corresponding to $v$ itself and the index
$j-1$. A nonzero value of $F_{0,K}$ at that pair already places $v$
in $l.W_t$. Its value there is $\mathrm{FK}(s,D,j,p)$.

Moreover $\mu_{\rm At}(v)=\mu_0$ forces $v\notin l.W_t$.
On $l.W_t$ the row identifies the character of that tuple with
the prescribed twist $\xi_{\rm Tw}(v)$, whereas $\xi(\mu_0)=1$.
The twist is nontrivial: it is the product of one of the two
distinguished characters and an allowed twisting character.
At a prime above the distinguished character's defining prime
the latter twisting factor is unramified by its exclusions, and
the distinguished factor has nontrivial inertia. Thus their product
cannot be one. Off $l.W_t$ the row's untwisted formulas give the
two required zeros $I_j(\mu_0,v)=j_w(\mu_0,j,v)=0$.

Consequently any supplied division record $Y$ satisfying the
row Hermitian-lattice comparisons $\mathrm{MjAre}$ and
$Y.\kappa_R(\sigma)=\kappa_{\rm sharp}(D)$ gives the core's
$\mathrm{RowLay}$. Indeed the sharp rational-roots condition
comes from the lift reading's rational-roots identity by the
sharp-row equivalence, the base rational-roots condition is its
base identity, and the four remaining row identities for
$\kappa,\lambda,b_u$ and occurrence come from the ordinary reading.
The two zeros prove its untwisted field. The row's index comparison
applied to $Y$ and the displayed $\kappa_R$ formula finish it.
% END P78C4

% BEGIN P78C5
\paragraph{P78C5. A built core retaining the kernel certificate.}
Fix a pinned fourth-stage record as in P78C3 in a context with
its original over-base and forced-prime conditions. Suppose the
division and theta constructors give their specified outputs on
this level. Thus one has $Y$, its division tie, its row metric and
$\kappa_R$ comparisons and its object-slot witness for any object
with these very division and row data; one also has $\theta$ with
its theta tie to $l.n$. Use the following geometric inputs at
their original scopes: a CM gauge on the host, the column-layer
identities for the stored split column, and the lattice-layer and
kernel construction with its certificate. The last is evaluated
\emph{after} taking $Y$, using its $Q_p$ tie. Its outputs are
one $G_y$, the kernel family for the stored comparisons and twists,
the lattice-layer identities, a kernel witness $k$ over $Y$, and
place data, matrices and a compatibility witness with
$k.E=\mathrm{kerEntC}(\mathrm{Pl},k.O,B_s,H)$.

Form the object $O$ using the stored column's group and comparison
data, its lower level and error, the column's twists, $l.nch,l.xd$,
$l.p_0$, the just chosen $G_y$, kernel, $Y$, $\theta$, gauge and
the two forms of P78C3. The division constructor's universal
object implication applies now with $O.X=Y$, $O.p_0=l.p_0$ and
the identical row packet. This order avoids prescribing an
object-dependent slot before $O$ exists.

Take the original rows, column character, set $l.W_t$, its tuple
comparison and the over-base and forced-prime conditions.
P78C4 gives the forced-content and row identities. The column
and lattice-layer inputs give the two layer comparisons. The group
stage gives the numerical group identity, and the ordinary reading
gives stability of the source module. The theta witness gives its
intertwining relation. The scale, archimedean signatures, real
comparison, $\eta$-factor and finite-type characters are those
of $l$; the choice $l.n$ gives its character tie and its local
parameters on $l.W_t$. This is a core $c$ with object $O$.
Pair it with the pin, coherent twists, original row, character
and spectral slot witnesses, and the division, theta and object
ties just constructed. The dependent pair of comparisons and
kernel required in the built witness is equal by construction
to the one on which $k$ was chosen. Thus the same $k.E$ proves
$\mathrm{EntC}$ for this built witness. This gives
$\mathrm{BuiltE}(\mathcal S,K,w,Z,c)$ at the specified level,
retaining the division, theta, column, gauge and kernel inputs.
% END P78C5

% BEGIN P78C6
\paragraph{P78C6. Converting the strengthened core into a good core for the original slots.}
For the second pass use the extension $N_g=\mathrm{NEog}$ and the
first-pass extension $N=\mathrm{NEo}$ at the same depth. The
second is a factor of the first:
\[
  N_g(\ldots,n,xd,c_0)
       = N(\ldots,n,xd,c_0)\times
            \mathrm{PLift}(\mathrm{GoodO}(D,nch,xd)).
\]
Here $n$ is the tie for the chosen $nch$, and $c_0$ is the
character-choice input. If a built core for the strengthened
slots has the kernel certificate, its very same $c$ is a good
core for the original slots, and $\mathrm{TwAllAt}$ holds for
$c.O.nch,c.O.xd$. To prove this, retain both the common part of
its character choice and the first projection in the displayed
product. They give the weaker $N$ witness on the identical $nch$.
Every other position, including the kernel family and its proof
$\mathrm{EntC}$, then has exactly the same type and value. The
second projection gives $\mathrm{GoodO}$, including the separation
of the prime selections when the two twist indices differ.

The common character choice has conjugate choice packets and its
side bit switches at distinct conjugate twisted places. These are
the two conjugation clauses of $\mathrm{TwAllAt}$. The twist field
of $\mathrm{GoodO}$ supplies
$I=(E,N_0,S)$ with $E=\mathrm{twEOf}(c.O.xd)$ and
$\mathrm{GoodAllAt}(D,c.O.nch,I)$. Substitute that equality
and keep $N_0,S$. This supplies the other clause of
$\mathrm{TwAllAt}$ with its original finite selection sets.
The conversion is pointwise, before taking any scale threshold.
% END P78C6

% BEGIN P78C7
\paragraph{P78C7. The order of the two passes through the large-field condition.}
Here is a sufficient construction of $\mathrm{MemberE}$ for the
two extensions of P78C6, specifying precisely the inputs retained.
Assume the original $\mathrm{BaseCond}(KPx)$ and eventual existence
of a pinned column over untwisted data; the first-pass step from
$\mathrm{Start}$ to a third-stage record; the progression from a
pinned third to fourth stage for each extension; and the compatible
division, theta, column, gauge and kernel inputs of P78C5 in both
passes. Assume also the export implication for every first-pass
built core, and the retwisting implication $\mathrm{Pass}$ from
any such exported core to a new starting column with a third-stage
record for the strengthened slots. Then $\mathrm{MemberE}$ holds
for the original slots.

Indeed an untwisted datum satisfies support by its untwisted
alternative and satisfies coherence: its prime and character
assignments agree on equivalent places and the required root has
order $p^0=1$. Retain the pin when applying the two first-pass
steps. P78C3 and P78C5 give the level, core and certificate
$\mathrm{BuiltE}$ there. Present its actual column and level as
the staged core witness. This proves the first pass. On the
enlarged base the same argument from the supplied strengthened
third-stage record gives a built core there. P78C6 makes it
simultaneously built, twist-compatible and good for the weaker
slots, with a single staged presentation.

For completeness, the eventual requirements of the second
construction have to be met on the \emph{output} of retwisting.
Intersect the context conditions for the two implications.
In the second construction impose the forced-prime clause by
including the two $g$ primes, the discriminant prime and the
integers up to $\lfloor K.s.X''(n)\rfloor$ in the finite prime
request. Write its four remaining finite requests as
$(S_1,P_1,A_0,S_{20})$, for small-field elements, primes,
avoidance and final field elements. Impose $S_1,P_1$ on the
input of retwisting. In its universal test, when the caller
asks for $A,S_2$, request the pointwise unions
$A\cup A_0$ and $S_2\cup S_{20}$ there. The latter remains
finite and algebraic. Use the original prehost, frame and small
fields in $A_0$; in $S_{20}$ allow the final character family
to vary with the enlargement, exactly as in the request.
The enlargement keeps the prehost and frame, both specified small
fields and the set of allowed primes. The new base consequently
meets $S_1,P_1$ and still has $KPx$ by the enlargement clause of
$\mathrm{BaseCond}$. Its avoidance and large coefficient field
meet the second construction's requests by inclusion in the
two unions. They also meet the caller's requests. Thus the
strengthened core constructed on that very base proves the
required $\mathrm{SecondG}$. Combine this with $\mathrm{BaseCond}$,
the first pass, and the assumed export. This proves the claimed
implication without assuming $\mathrm{AmbEx}$ or compatibility
between two separately chosen cores.

Uniform production of the input stages, their division, theta
and kernel data, and the retwisting and export hypotheses is
still needed to apply this implication. A large bundle at the
prescribed reciprocity record is the separate assertion in S77C5.
% END P78C7

% BEGIN P78C8
\paragraph{P78C8. The divisor retains the actual row lattices and one common type label.}
Fix a level-four input of P78C5, its stored row reading $Z$ and its
distinguished coefficient embedding $e_0$. For a finite type $\sigma$,
a label $c$ is admissible if, for every $i<m$, the $\sigma$-isotypic
subspace in row $i+1$ lies in the complex span of
$jW_i(W_{i,c})$. The set of all labels is nonempty. Indeed, an empty
label set would eliminate all good base eigensystems. The same-types
condition on the base then makes the good coordinate space zero, contrary
to positivity on the base side of these rows. Choose $c_\sigma$ admissible
when one exists, and otherwise choose any label. This choice precedes the
choice of row and embedding.

At $e_0$ put $\log\lambda(j,\sigma,e_0)=
\log(\lambda_P(\mu_0,j,\sigma))$. At $e\ne e_0$ put
\[
 \log\lambda(j,\sigma,e)=
 \begin{cases}
 \log(\lambda_{j-1,\mu_0,c_\sigma,[e]})&
 j-1<m\text{ and the isotypic space is nonzero},\\
 0&\text{otherwise}.
 \end{cases}
\]
Here natural subtraction is truncated, and $[e]$ denotes the associated
infinite place. At $j=i+1$ this is precisely each of the two clauses of
$\mathrm{LamAre}$ with the same $c_\sigma$ throughout. Set every
$\kappa_R(\sigma)=\kappa_\sharp$. For $1\le j=i+1\le m$ take the
Hermitian lattice of the stored $M_i,XW_i$ and its given semilinear map
$jW_i$. At other indices take the standard unit lattice of the complex
dimension of the row and map its basis to an orthonormal basis. The images
form a complex basis in either case. In range the integral membership
test, lattice inclusion, and metric equality are the existing pins and
goodness of that very row. Identifying coordinates by the identity on
$M_i$ proves $\mathrm{MjAre}$. These constructions supply the row part of
the divisor without requiring local integral values.
% END P78C8

% BEGIN P78C9
\paragraph{P78C9. The ramified places and their selected primes.}
Let $E/F$ be the input CM extension, $\xi_\mu$ the tied character of
the actual tuple, and $P_v$ the prime chosen by the column above $v$.
The ramification pins follow from the input $\mathrm{MuIsW}$. Every
$\xi_\mu$ belongs to its specified twist family, so a ramified prime has
prime absolute norm $p$, satisfies $p\notin P^2$, and avoids the two
host primes. The character of $\mu(w)$ is unramified above the residue
characteristic of $w$. On the twisting set this is the prescribed local
condition on the twist, and off that set it is the trivial character.

Let $V_\mu$ be the set of primes $v$ of $F$ below a ramified prime of
$\xi_\mu$. It is finite by finiteness of the conductor. If $w$ above
$v$ is ramified, so is $P_v$. For distinct $w,P_v$ these exhaust the
primes over $v$, since $[E:F]=2$. If the character at $P_v$ were trivial
on units, choose any unit $u$ at $w$. Its residue is that of an ordinary
integer $a$ prime to $p$, since $w$ has prime residue cardinality.
Tameness identifies the value at $u$ with that at $a$. The integer $a$
is a unit at both primes. Evaluating the restriction of the character
to the base idele which equals $a$ at $v$ and $1$ elsewhere gives
$\xi_{\mu,w}(a)\xi_{\mu,P_v}(a)=1$. That restriction is trivial.
This would make $\xi_{\mu,w}$ trivial on all units, a contradiction.

The pins now apply at $P_v$ itself. Its residue field has $p$ elements,
and $p\in P_v\setminus P_v^2$ gives $\operatorname{ord}_{P_v}(p)=1$.
The residue field at $v$ embeds in this prime field and hence also has
$p$ elements. Thus $q_v=Nv=NP_v=p$ and $P_v$ has the required degree
one. The rational support is exactly $\{NP_v:v\in V_\mu\}$. For either
inclusion use a ramified prime and its contraction to $F$; all primes
above that contraction have the same rational residue characteristic.
This is the column's prescribed $\mathrm{QpIs}$ set under the given tie
of the column and tuple characters. It is empty for the base tuple and
has the two host-prime exclusions and the $\mu(w)$ exclusion above.
% END P78C9

% BEGIN P78C10
\paragraph{P78C10. The first local doubling value.}
Over the completion at a degree-one $P$, write $O$ for the integers,
$\varpi$ for the chosen uniformizer and $q=NP$. Let $\delta$ be the
ramified tame unitary character supplied at $P$, and $\nu$ any of the
prescribed unramified unitary characters. For rank $r\ge0$ and block
length $t\ge1$ the single-block local step, with constant character
family $\nu$, has the value
\[
             q^{-tr}\prod_{h=1}^{r}(1-q^{-h}).
\]
In the index-zero step the second Levi block has size zero. The trivial
one-dimensional representation and the constant vector function $f=1$
satisfy the type conditions, since the right character and the modulus
in these conditions are trivial. Its spherical coefficient is $1$.
Normalize Haar measure by $\mu(K)=1$, $K=\mathrm{GL}_r(O)$, so also
$\mathrm{mcoef}(f)=1$. In row-block coordinates $x=(A;B)$ the test
function equals $\delta(\det A)^{-1}$ for $A\in K$ and
$B\in\varpi M_{t,r}(O)$, and vanishes elsewhere. The two functions in
the doubling integral can overlap only if $A,Ag\in K$, hence $g\in K$.
For $g\in K$ right translation preserves the support and contributes
$\delta(\det g)^{-1}$. This cancels the $\delta$ part of
$(\nu\delta)(\det g)$; both $\nu(\det g)$ and the absolute-value
factor equal $1$. Unitarity removes the remaining paired phases at $A$.
All integrations here have compact support.

Give each additive block of integral matrices volume one and concatenate
the two row blocks to obtain the required additive measure. The lower
block then has volume $q^{-tr}$. The admissible upper blocks have
volume $|\mathrm{GL}_r(\mathbb F_q)|/q^{r^2}$, since each residue
class of matrices has equal measure. Successively choosing independent
columns gives the displayed product, also for $r=0$ with the empty
product convention. The group integral contributes $\mu(K)=1$.
This evaluates the actual first quotient observable in the divisor
record, including its Haar normalizations.
% END P78C10

% BEGIN P78C11
\paragraph{P78C11. Witt embeddings that read the prescribed coefficient place.}
Fix a rational prime $p$ and a place $\lambda\mid p$ of the coefficient
number field $L$. Write $W=W(\overline{\mathbb F}_p)$ and $B=\mathrm{Frac}(W)$.
There exists an embedding $j:B\hookrightarrow\mathbb C$ and a valuation
ring $V$ of $\mathbb C$ whose pullbacks under $j$ and the distinguished
embedding $e_0:L\hookrightarrow\mathbb C$ are $W$ and $O_{L,\lambda}$
respectively. Indeed, first embed $B$ abstractly in $\mathbb C$.
The set of its elements has cardinality at most the continuum, so a
transcendence basis embeds into one of $\mathbb C/\mathbb Q$, and the
algebraic part embeds since $\mathbb C$ is algebraically closed. Extend
the valuation ring $W$ along this embedding by valuation prolongation.
In the resulting ring $V_0$, $p^{-1}$ is absent; it is already absent
from $W$. On $\overline{\mathbb Q}$ the induced valuation and an
extension of $\lambda$ both lie over the $p$-adic valuation. They
are conjugate: on each finite normal subextension the Galois group is
transitive on primes over $p$, and compactness of the absolute Galois
group gives a single conjugator for all finite subextensions. Thus
there is an embedding $e:L\hookrightarrow\mathbb C$ pulling $V_0$
back to $O_{L,\lambda}$. An automorphism of $\mathbb C$ takes $e$ to
$e_0$, by extending the isomorphism on algebraic closures and then on
transcendence bases. Transport both the initial embedding and $V_0$
by that automorphism.

In particular, whenever $j(y)=e_0(z)$ one has
$y\in pW$ if and only if $z$ belongs to the maximal ideal of
$O_{L,\lambda}$. For a nonzero element both tests say its inverse is
outside the respective valuation ring; zero passes both tests. A Witt
vector belongs to $pW$ precisely when its first coefficient vanishes,
because the residue map $W\to\overline{\mathbb F}_p$ is the first
coefficient and has kernel $pW$. Apply this on each common reading of
an integral host element and choose the embeddings independently for
all $(p,\lambda)$. This proves $\mathrm{ComplexEmbeddingsCompatible}$
with the fixed coefficient embedding throughout.
% END P78C11

% BEGIN P78C12
\paragraph{P78C12. The prescribed cyclic spherical representation at every block length.}
Write $k=E_P$, $O=O_k$, $|\varpi|=q^{-1}$, $c=jt$, $G=\mathrm{GL}_c(k)$
and $K=\mathrm{GL}_c(O)$. Let $P_0$ be the upper parabolic with $j$
consecutive blocks of size $t\ge1$. Write $A_b$ for its diagonal blocks
with $0\le b<j$. Under conjugation by a block diagonal element the
upper off-diagonal block $(b,d)$ changes by $X\mapsto A_bXA_d^{-1}$.
Its additive measure scales by
$|\det A_b|^t|\det A_d|^{-t}$. The product is the parabolic modulus
$D(p)=\prod_b|\det A_b|^{c-t-2bt}$. For the input unitary unramified
characters put
\[
 \psi(p)=D(p)^{1/2}\prod_b\nu_{b+1}(\det A_b).
\]
This character equals $1$ on $P_0\cap K$. Every element is in $P_0K$:
the intersections of an $O$-lattice with the successive spaces of a
flag are saturated, and the resulting free quotients over the DVR allow
one to choose an adapted integral basis. Equivalently $K$ acts
transitively on the partial flags and $G=P_0K$ by inversion.
Thus $\phi(pk)=\psi(p)$ is well defined, right $K$-invariant, and takes
the value $1$ on $K$.

Give the complex span $V$ of the right translates of $\phi$ the pairing
$\langle u,v\rangle=\int_K\overline{u(k)}v(k)\,dk$ with $K$ of
volume one. Each vector is fixed by a finite intersection of open
conjugates of $K$, hence is locally constant. The pairing is positive
definite, since a nonzero value on $K$ persists on an open subset of
positive measure and restriction to $K$ determines a function by
$P_0$-covariance.

Here is the invariance under right translations. Give $P_0$ left Haar
measure with $P_0\cap K$ of volume one. For compactly supported locally
constant $h$ on $G$,
\[
       \int_G h(x)\,dx=\int_K\int_{P_0} h(pk)\,dp\,dk.
\]
Indeed, both measures are left $P_0$-invariant, right $K$-invariant,
and give $K$ mass one. Every coset in the discrete $G/K$ has a
representative in $P_0$ and on each coset right $K$-invariance
determines the measure. For $F=\overline u v$ one has
$F(px)=D(p)F(x)$. Taking $h=1_K F$ gives
$F(x)=\int_{P_0}h(px)\,dp$: equality holds on $K$ by normalization
and extends by this same covariance. For the latter, right translation
in left Haar integration on $P_0$ contributes $D(p)$, as follows from
the conjugation factors on the upper unipotent group computed above
and unimodularity of the diagonal blocks. Apply the displayed
disintegration to $h(\cdot\,g)$. Haar measure on $G$ is right
invariant as well, as seen from $|\det x|^{-c}\,dx$ in matrix
coordinates. This proves $\int_K F(kg)\,dk=\int_K F(k)\,dk$ and makes
the right translations isometries.

The fixed line is $\mathbb C\phi$ since a $K$-fixed vector is
determined on $P_0K$ by its value at $1$. For a nonzero invariant
subspace $U$, a nonzero member is nonorthogonal to some translate of
$\phi$ by cyclicity. Translate that member back to $w\in U$ with
$\langle\phi,w\rangle\ne0$. Its $K$-orbit is finite by smoothness and
compactness. The sum of the distinct orbit elements is fixed and its
pairing with $\phi$ is the orbit cardinality times
$\langle\phi,w\rangle$. Thus $U$ contains $\phi$ and then all of $V$.
This proves algebraic irreducibility on the actual smooth space.

Let $\mathrm{GL}_r(k)$ act trivially on this representation. The unit
vector $\phi$ is spherical and its coefficient at $(1,g)$ is
$\int_K\phi(kg)\,dk$. To identify the requested parameter at diagonal
entry $x=bt+s$, $0\le s<t$, multiply the defining power of
$\mathrm{satC}$ by the normalized Borel power. The exponent becomes
\[
       s-(t-1)/2+(c-1)/2-x=(c-t)/2-bt.
\]
The character is $\nu_{b+1}$. Multiplication of these factors across
the block gives exactly $\psi$ for an upper triangular matrix, with
right $K$-invariance as above. This proves every clause of
$\mathrm{IsUnrM}$ at the specified parameters. For $j=0$ use the
trivial group and its unit vector as in P78C10.
% END P78C12

% BEGIN P78C13
\paragraph{P78C13. A normalized type vector for every required local step.}
Continue with the model $(\rho,v_0)$ of P78C12, $\|v_0\|=1$ and $c=jt$.
In $\mathrm{GL}_{r+c}(k)$ let $P_1$ have blocks $(r,c)$ and set
$K_1=\mathrm{GL}_{r+c}(O)$. Let $J\subset K_1$ be the preimage
of the upper block parabolic modulo $\varpi$. The character
$\lambda_J(h)=\delta(\det h_{22})$ is multiplicative on $J$:
the lower-right blocks multiply modulo $\varpi$, their determinants
are units, and $\delta$ is trivial on $1+\varpi O$. For $p\in P_1$
with diagonal blocks $a,d$, set
\[
 f(ph)= |\det a|^{c/2}|\det d|^{-r/2}
        \delta(\det d)\lambda_J(h)\rho(a,d)v_0
        \quad(h\in J),
\]
and put $f=0$ off $P_1J$. On $P_1\cap J$ both diagonal blocks are
integral invertible matrices. They fix $v_0$, have determinant norms
one and give the same character on the two sides of this formula.
Any two decompositions differ by such an element, proving
well-definedness. Left parabolic covariance and right $J$-equivariance
then give the defining type. On $K_1$ the norm of $f$ equals
$1_{K_1\cap P_1J}$ by unitarity and the same integral block test.

Normalize left Haar measure by $\mu(J)=1$. The open set
$K_1\cap P_1J$ contains $J$ and has finite volume $C$ since $K_1$
is compact. Thus $C\ge1$ and $f'=C^{-1/2}f$ has
$\mathrm{mcoef}(f')(1)=\int_{K_1}\|f'(k)\|^2\,d\mu=1$.
Normalize additive matrix Haar measure to give the integral
$(r+(j+1)t)\times(r+jt)$ matrices volume one. Take $n$ to be the
specified $\mathrm{finrank}_{\mathbb C}$ of this type space and take
$z$ equal to the local doubling integral of the source test functions
with this normalized coefficient and $\nu_{j+1}\delta$. These choices
prove $\mathrm{TypeDimAt}(n)$ and $\mathrm{StepAt}(z)$ as defined.
For $j>0$ the construction requires existence of these stored values;
an evaluation of the integral or a formula for the dimension is a
separate assertion.
% END P78C13

% BEGIN P78C14
\paragraph{P78C14. The divisor producer on the two passes.}
For every original level-four input of P78C5, both native slot choices
admit the very divisor required by $\mathrm{DivEx}$. Begin with the
row record of P78C8. For each $\mu$ replace its prime set and its local
indices by the rational support and the set $V_\mu$ of P78C9.
Use $q_v=Nv$ there. At the chosen $P_v$ set
$\delta=(\xi_\mu^{-1})_{P_v}$, which is ramified, tame and unitary.
For $1\le i\le m$ take $\nu_i=((C.\mu_0(i-1))^{-1})_{P_v}$
from the frame $C$, and take $\nu_i=1$ otherwise. These characters
are unitary and unramified since the two host primes are excluded.
Thus $\nu_{i}\delta$ is the inverse local character of the
corresponding tuple entry for $1\le i\le m$. Store $\nu_i(\varpi_{P_v})$,
nonzero and of absolute value one. Use P78C13 to choose a dimension $n_j$ and a value $z_j$ for
each step and set $Z_Q(0)=0$, $Z_Q(j+1)=z_j$. For the separate first
quotient observable use P78C10, applicable with each requested
unramified character. The local tie follows from the defining choices,
and the identification of rational primes and $q_v$ is P78C9.

First construct the divisor for $\mathrm{OX}(\mathrm{anyGeometry})$.
The geometric model is explicitly permitted by $\mathrm{GeoTiedP}$. Take the supplied small coefficient field with
its inclusion in $L$ and its own ring of integers. Use the flag family
of the source with singleton level, level primes the two host primes,
and the Witt embeddings of P78C11. Its finite-place index at $p$
consists of the places of $L$ above $p$ with their original normalized
weights. At each row the component index is a singleton, the map from
the row lattice to the component's sections is the zero linear map,
and the allowed sections contain zero. The local model and curve
indices are singletons and their degree is zero. In this family the
tests attached to a flag use its rank and index. The stored flag on
a test with the prescribed row has these same two indices, so it gives
the required equality of flag objects and local models. The
component, boundary and weight-line objects are those of that same
family by construction. This also proves invariance under
isomorphisms preserving the rank and index. The level membership
predicate is identically true, the vanishing condition follows from
the zero section maps, and the place compatibility is P78C11. These
are the component, local, cusp, invariance, and level clauses of the
specified geometric tie.

Graft these geometric fields onto the arithmetic record, retaining
its row lattices, row maps, logarithms, constants, prime support and
local observables. The component section maps just specified take
the retained row lattice as domain, so the graft is defined on the
original input. P78C8 gives $\mathrm{MjAre}$, every $\kappa_R$ and
$\mathrm{LamAre}$ after grafting. If the divisor is tested on an
equal row side and an equal transported reading, substitute those
equalities to obtain the same $\mathrm{LamAre}$ witness. The second
part of the native slot test is then satisfied, and
$\mathrm{anyGeometry}$ takes the value true on this witness. For the
actual slots $\mathrm{OX}(\mathrm{geoCl}(\mathrm{anyGeometry}))$,
use the base constructor of $\mathrm{geoCl}$ on this same witness.
The two slot choices carry identical input data in this conversion. The
construction is uniform in the original level-four input; it applies
to both passes of P78C7. Its divisor hypothesis is therefore discharged.
The construction establishes the geometric predicate that these slots
ask for; stronger moduli interpretations require their own hypotheses.
% END P78C14
